\Article{20}
Ist die Anzahl der Bedingungsgleichungen sehr groß, so kann die unmittelbare Elimination zur Bestimmung der Korrelate \(A,B,C\) etc. so langwierig werden, dass die Geduld des Rechnenden dafür nicht ausreicht. Dann kann es häufig zweckmäßig sein, mithilfe des Lehrsatzes im vorhergehenden Artikel durch sukzessive Annäherungen eine vollständige Ausgleichung zu ermitteln. Man teile die Bedingungsgleichungen in zwei oder mehrere Klassen ein und bestimme zunächst eine Ausgleichung, durch die den Gleichungen der ersten Klasse genügt wird, während die übrigen außer Betracht bleiben. Danach behandle man die durch diese Ausgleichung geänderten Beobachtungen so, dass nur die Gleichungen der zweiten Klasse berücksichtigt werden. Allgemein gesprochen wird die Anwendung des zweiten Ausgleichungssystems die Übereinstimmung mit den Gleichungen der ersten Klasse stören; wenn also nur zwei Klassen gebildet wurden, kehren wir zu den Gleichungen der ersten Klasse zurück und ermitteln ein drittes System, das ihnen genügt. Dann unterziehen wir die dreimal korrigierten Beobachtungen einer vierten Ausgleichung,
\SourceLeaf{084}
bei der nur die Gleichungen der zweiten Klasse berücksichtigt werden. Auf diese Weise erhalten wir abwechselnd, bald die erste, bald die zweite Klasse berücksichtigend, stetig abnehmende Ausgleichungen; und wenn die Einteilung sachkundig vorgenommen wurde, werden wir nach wenigen Iterationen stabile Werte erreichen. Sind mehr als zwei Klassen gebildet worden, verhält es sich ähnlich: Die einzelnen Klassen gehen nacheinander in die Rechnung ein, nach der letzten geht wieder die erste in die Rechnung ein und so weiter. Doch soll es hier genügen, dieses Verfahren angedeutet zu haben; seine Wirksamkeit wird sicherlich in hohem Maße von einer sachkundigen Anwendung abhängen.

\Article{21}
Es bleibt noch, den Beweis des in Art. 8 vorausgesetzten Lemmas zu ergänzen; der Deutlichkeit halber werden wir dabei jedoch andere, diesem Gegenstand besser angepasste Bezeichnungen verwenden.

Seien also \(x^0,x',x'',x'''\) etc. Unbestimmte, und es werde angenommen, dass sich aus den Gleichungen
\[
\begin{aligned}
n^{00}x^0+n^{01}x'+n^{02}x''+n^{03}x'''+\text{etc.}&=X^0\\
n^{10}x^0+n^{11}x'+n^{12}x''+n^{13}x'''+\text{etc.}&=X'\\
n^{20}x^0+n^{21}x'+n^{22}x''+n^{23}x'''+\text{etc.}&=X''\\
n^{30}x^0+n^{31}x'+n^{32}x''+n^{33}x'''+\text{etc.}&=X'''\\
\text{etc.}&
\end{aligned}
\]
durch Elimination die folgenden Gleichungen ergeben:
\[
\begin{aligned}
N^{00}X^0+N^{01}X'+N^{02}X''+N^{03}X'''+\text{etc.}&=x^0\\
N^{10}X^0+N^{11}X'+N^{12}X''+N^{13}X'''+\text{etc.}&=x'\\
N^{20}X^0+N^{21}X'+N^{22}X''+N^{23}X'''+\text{etc.}&=x''\\
N^{30}X^0+N^{31}X'+N^{32}X''+N^{33}X'''+\text{etc.}&=x'''\\
\text{etc.}&
\end{aligned}
\]

Setzen wir nun in die erste und zweite Gleichung des zweiten Systems die Werte der Größen \(X,X',X'',X'''\) etc. aus dem ersten System ein, so erhalten wir
\[
\begin{aligned}
x^0={}&N^{00}(n^{00}x^0+n^{01}x'+n^{02}x''+n^{03}x'''+\text{etc.})\\
&+N^{01}(n^{10}x^0+n^{11}x'+n^{12}x''+n^{13}x'''+\text{etc.})\\
&+N^{02}(n^{20}x^0+n^{21}x'+n^{22}x''+n^{23}x'''+\text{etc.})\\
&+N^{03}(n^{30}x^0+n^{31}x'+n^{32}x''+n^{33}x'''+\text{etc.})\quad\text{etc.}
\end{aligned}
\]
\SourceLeaf{085}
und außerdem
\[
\begin{aligned}
x'={}&N^{10}(n^{00}x^0+n^{01}x'+n^{02}x''+n^{03}x'''+\text{etc.})\\
&+N^{11}(n^{10}x^0+n^{11}x'+n^{12}x''+n^{13}x'''+\text{etc.})\\
&+N^{12}(n^{20}x^0+n^{21}x'+n^{22}x''+n^{23}x'''+\text{etc.})\\
&+N^{13}(n^{30}x^0+n^{31}x'+n^{32}x''+n^{33}x'''+\text{etc.})\quad\text{etc.}
\end{aligned}
\]

Da jede dieser Gleichungen offensichtlich eine identische Gleichung sein muss, können wir sowohl in der ersten als auch in der zweiten für \(x^0,x',x'',x'''\) etc. beliebige bestimmte Werte einsetzen. In die erste setzen wir
\[
x^0=N^{10},\quad x'=N^{11},\quad x''=N^{12},\quad x'''=N^{13}\quad\text{etc.}
\]
in die zweite hingegen
\[
x^0=N^{00},\quad x'=N^{01},\quad x''=N^{02},\quad x'''=N^{03}\quad\text{etc.}
\]
Nach dieser Substitution ergibt die Subtraktion
\[
\begin{aligned}
N^{10}-N^{01}={}&(N^{00}N^{11}-N^{10}N^{01})(n^{01}-n^{10})\\
&+(N^{00}N^{12}-N^{10}N^{02})(n^{02}-n^{20})\\
&+(N^{00}N^{13}-N^{10}N^{03})(n^{03}-n^{30})\\
&+\text{etc.}\\
&+(N^{01}N^{12}-N^{11}N^{02})(n^{12}-n^{21})\\
&+(N^{01}N^{13}-N^{11}N^{03})(n^{13}-n^{31})\\
&+\text{etc.}\\
&+(N^{02}N^{13}-N^{12}N^{03})(n^{23}-n^{32})\\
&+\text{etc. etc.}
\end{aligned}
\]
Diese Gleichung kann auch in der Form
\[
N^{10}-N^{01}=\Sigma(N^{0\alpha}N^{1\beta}-N^{1\alpha}N^{0\beta})(n^{\alpha\beta}-n^{\beta\alpha})
\]
dargestellt werden, wobei \(\alpha\beta\) alle Kombinationen ungleicher Indizes bezeichnen.

Hieraus folgt, dass, wenn
\[
\begin{aligned}
&n^{01}=n^{10},\quad n^{02}=n^{20},\quad n^{03}=n^{30}\\
&n^{12}=n^{21},\quad n^{13}=n^{31},\quad n^{23}=n^{32},\quad\text{etc.}
\end{aligned}
\]
oder allgemeiner \(n^{\alpha\beta}=n^{\beta\alpha}\) gilt, auch
\[
N^{10}=N^{01}
\]
\SourceLeaf{086}
gelten wird. Da die Reihenfolge der Unbestimmten in den vorgelegten Gleichungen beliebig ist, gilt unter jener Voraussetzung offenbar allgemein
\[
N^{\alpha\beta}=N^{\beta\alpha}
\]
